
\documentclass[11pt]{article}

\usepackage[a4paper,margin=25mm]{geometry}
\usepackage[T1]{fontenc}
\usepackage{lmodern}
\usepackage{amsmath,amssymb,amsthm,mathtools}
\usepackage{microtype}
\usepackage{enumitem}

\setlength{\parindent}{0pt}
\setlength{\parskip}{0.45em}
\allowdisplaybreaks

\newtheorem{lemma}{Lemma}
\newcommand{\norm}[1]{\left\lVert #1\right\rVert}
\newcommand{\ct}{\mathrm{ct}}
\newcommand{\ctadd}{\mathrm{ct}_{\mathrm{add}}}
\newcommand{\ctmult}{\mathrm{ct}_{\mathrm{mult}}}
\newcommand{\ctrelin}{\mathrm{ct}_{\mathrm{relin}}}
\newcommand{\ctpmult}{\mathrm{ct}_{\mathrm{pmult}}}
\newcommand{\ctpadd}{\mathrm{ct}_{\mathrm{padd}}}
\newcommand{\ctneg}{\mathrm{ct}_{\mathrm{neg}}}
\newcommand{\ctsub}{\mathrm{ct}_{\mathrm{sub}}}
\newcommand{\ctpsub}{\mathrm{ct}_{\mathrm{psub}}}
\newcommand{\vadd}{v_{\mathrm{add}}}
\newcommand{\vmult}{v_{\mathrm{mult}}}
\newcommand{\vrelin}{v_{\mathrm{relin}}}
\newcommand{\vpmult}{v_{\mathrm{pmult}}}
\newcommand{\vpadd}{v_{\mathrm{padd}}}
\newcommand{\vneg}{v_{\mathrm{neg}}}
\newcommand{\vsub}{v_{\mathrm{sub}}}
\newcommand{\vpsub}{v_{\mathrm{psub}}}
\newcommand{\N}{\mathbb{N}}

\title{SEAL 2.3.1 Manual: Appendix on Strict Noise Bounds}
\author{Transcription of Appendix, Lemmas 9--17}
\date{}

\begin{document}
\maketitle

\textbf{Source.} Kim Laine, \emph{Simple Encrypted Arithmetic Library (SEAL) 2.3.1},
Appendix, printed pages 27--34 (PDF pages 27--34).

\textbf{Scope.} This file transcribes the appendix lemmas and their accompanying
proofs or explanations concerning strict bounds on invariant noise. The
manual's notation and lemma numbering are retained. No heuristic bounds from
Section 6 are substituted.

\section*{Initial Noise}

\begin{lemma}[Lemma 9]
Let $\ct=(c_0,c_1)$ be a fresh encryption of a message $m\in R_t$.
The noise $v$ in $\ct$ satisfies
\[
  \norm{v}
  \leq \frac{r_t(q)}{q}\norm{m}
       +\frac{tB}{q}(2n+1).
\]
\end{lemma}

\textbf{Proof.}
Let $\ct=(c_0,c_1)$ be an encryption of $m$ under the public key
$\mathrm{pk}=(p_0,p_1)=([-(as+e)]_q,a)$. Then, for some polynomials
$k_0,k_1,k$,
\begin{align*}
\frac{t}{q}(c_0+c_1s)
&=\frac{t}{q}
  \bigl(\Delta m+p_0u+e_0+k_0q+p_1us+e_1s+k_1qs\bigr)\\
&=m+\frac{t}{q}
  \left(\frac{-r_t(q)m}{t}+p_0u+e_0+p_1us+e_1s\right)
  +t(k_0+k_1s)\\
&=m+\frac{t}{q}
  \left(\frac{-r_t(q)m}{t}+(-as-e+kq)u+e_0+aus+e_1s\right)
  +t(k_0+k_1s)\\
&=m+\frac{t}{q}
  \left(\frac{-r_t(q)m}{t}-eu+e_1+e_2s\right)
  +t(k_0+k_1s+ku).
\end{align*}
Thus the noise is
\[
  v=\frac{t}{q}
  \left(\frac{-r_t(q)}{t}m-eu+e_1+e_2s\right).
\]
To bound $\norm{v}$, use the fact that the error polynomials sampled from
$\chi$ have coefficients bounded by $B$, and that $\norm{s}=\norm{u}=1$.
Then
\[
  \norm{v}\leq \frac{r_t(q)}{q}\norm{m}
                  +\frac{tB}{q}(2n+1).
\]
\hfill$\square$

\section*{Addition}

\begin{lemma}[Lemma 10]
Let $\ct_1=(c_0,c_1,\ldots,c_j)$ and
$\ct_2=(d_0,d_1,\ldots,d_k)$ be two ciphertexts encrypting
$m_1,m_2\in R_t$, and having noises $v_1,v_2$, respectively.
Then the noise $\vadd$ in their sum $\ctadd$ is
\[
  \vadd=v_1+v_2,
\]
and satisfies
\[
  \norm{\vadd}\leq \norm{v_1}+\norm{v_2}.
\]
\end{lemma}

\textbf{Proof.}
By definition of homomorphic addition, $\ctadd$ encrypts $[m_1+m_2]_t$.
Let $[m_1+m_2]_t=m_1+m_2+at$ for some integer-coefficient polynomial $a$.
Suppose without loss of generality that $\max(j,k)=j$, so that
\[
  \ctadd=(c_0+d_0,\ldots,c_k+d_k,c_{k+1},\ldots,c_j).
\]
By definition of noise in $\ct_1$ and $\ct_2$,
\[
  \frac{t}{q}\ct_1(s)=m_1+v_1+a_1t,\qquad
  \frac{t}{q}\ct_2(s)=m_2+v_2+a_2t,
\]
for some integer-coefficient polynomials $a_1,a_2$. Therefore
\begin{align*}
\frac{t}{q}\ctadd(s)
&=\frac{t}{q}\ct_1(s)+\frac{t}{q}\ct_2(s)\\
&=m_1+v_1+a_1t+m_2+v_2+a_2t\\
&=[m_1+m_2]_t+(m_1+m_2-[m_1+m_2]_t)+v_1+v_2+(a_1+a_2)t\\
&=[m_1+m_2]_t+v_1+v_2+(a_1+a_2-a)t.
\end{align*}
Hence $\vadd=v_1+v_2$, and
\[
  \norm{\vadd}=\norm{v_1+v_2}
  \leq \norm{v_1}+\norm{v_2}.
\]
\hfill$\square$

\section*{Multiplication}

\begin{lemma}[Lemma 11]
Let $\ct_1=(x_0,\ldots,x_{j_1})$ be a ciphertext of size $j_1+1$
encrypting $m_1$ with noise $v_1$, and let
$\ct_2=(y_0,\ldots,y_{j_2})$ be a ciphertext of size $j_2+1$
encrypting $m_2$ with noise $v_2$. Let $N_{m_1}$ and $N_{m_2}$
be upper bounds on the number of non-zero terms in the polynomials
$m_1$ and $m_2$, respectively. Then the noise $\vmult$ in the product
$\ctmult$ satisfies
\begin{align*}
\norm{\vmult}\leq{}&
\left[
 (N_{m_1}+n)\norm{m_1}
 +\frac{nt}{2}\frac{n^{j_1+1}-1}{n-1}
\right]\norm{v_2}\\
&+\left[
 (N_{m_2}+n)\norm{m_2}
 +\frac{nt}{2}\frac{n^{j_2+1}-1}{n-1}
\right]\norm{v_1}\\
&+3n\norm{v_1}\norm{v_2}
 +\frac{t}{2q}\left(\frac{n^{j_1+j_2+1}-1}{n-1}\right).
\end{align*}
\end{lemma}

\textbf{Proof.}
By definition of homomorphic multiplication, the ciphertext
$\ctmult=(c_0,\ldots,c_{j_1+j_2})$ is such that for
$0\leq i\leq j_1+j_2$, for some polynomials $\epsilon_i$ with
coefficients in $(-\frac12,\frac12]$, and for some polynomials
$A_i$ with integer coefficients,
\[
 c_i=
 \left[\frac{t}{q}\left(\sum_{k+\ell=i}x_ky_\ell\right)\right]_q
 =\frac{t}{q}\left(\sum_{k+\ell=i}x_ky_\ell\right)+A_iq
 =\frac{t}{q}\left(\sum_{k+\ell=i}x_ky_\ell\right)+\epsilon_i+A_iq.
\]
Also, by definition $\ctmult$ encrypts $[m_1m_2]_t$, and
$[m_1m_2]_t=m_1m_2+at$ for some integer-coefficient polynomial $a$.
By definition of noise in $\ct_1$ and $\ct_2$, for some
integer-coefficient polynomials $a_1,a_2$,
\[
  \frac{t}{q}\ct_1(s)=m_1+v_1+a_1t,\qquad
  \frac{t}{q}\ct_2(s)=m_2+v_2+a_2t.
\]
We compute
\begin{align*}
\frac{t}{q}\ctmult(s)
&=\frac{t}{q}(c_0,\ldots,c_{j_1+j_2})(s)\\
&=\frac{t}{q}\left[
 \left(\frac{t}{q}(x_0y_0)+\epsilon_0+A_0q\right)
 +\cdots+
 \left(\frac{t}{q}(x_{j_1}y_{j_2})+\epsilon_{j_1+j_2}
       +A_{j_1+j_2}q\right)s^{j_1+j_2}
 \right]\\
&=\frac{t}{q}\cdot\frac{t}{q}
 \left[\sum_{i=0}^{j_1+j_2}
       \left(\sum_{k+\ell=i}x_ky_\ell\right)s^i\right]
 +\frac{t}{q}\sum_{i=0}^{j_1+j_2}\epsilon_i s^i
 +\left(\sum_{i=0}^{j_1+j_2}A_is^i\right)t\\
&=\frac{t}{q}\ct_1(s)\cdot\frac{t}{q}\ct_2(s)
 +\frac{t}{q}\sum_{i=0}^{j_1+j_2}\epsilon_i s^i
 +\left(\sum_{i=0}^{j_1+j_2}A_is^i\right)t\\
&=(m_1+v_1+a_1t)(m_2+v_2+a_2t)
 +\frac{t}{q}\sum_{i=0}^{j_1+j_2}\epsilon_i s^i
 +\left(\sum_{i=0}^{j_1+j_2}A_is^i\right)t\\
&=[m_1m_2]_t+m_1v_2+m_2v_1+v_1v_2
 +(v_1a_2+v_2a_1)t
 +\frac{t}{q}\sum_{i=0}^{j_1+j_2}\epsilon_i s^i\\
&\quad+\left(m_1a_2+m_2a_1+a_1a_2t
             +\sum_{i=0}^{j_1+j_2}A_is^i-a\right)t.
\end{align*}
In the last step we used $m_1m_2=[m_1m_2]_t-at$. Thus
\[
  \vmult=m_1v_2+m_2v_1+v_1v_2
       +(v_1a_2+v_2a_1)t
       +\frac{t}{q}\sum_{i=0}^{j_1+j_2}\epsilon_i s^i.
\]

To bound the new noise, first note that
\begin{equation}
 \frac{t}{q}\left\lVert\sum_{i=0}^{j_1+j_2}\epsilon_i s^i\right\rVert
 \leq \frac{t}{2q}
       \left(\frac{n^{j_1+j_2+1}-1}{n-1}\right).
 \tag{1}
\end{equation}
Next, write $a_it=\frac{t}{q}\ct_i(s)-m_i-v_i$, and note that
\begin{equation}
 \norm{a_it}\leq
 \frac{t}{2}\frac{n^{j_i+1}-1}{n-1}
 +\norm{m_i}+\norm{v_i}.
 \tag{2}
\end{equation}
Finally, using (1) and (2), we can bound the noise growth in multiplication:
\begin{align*}
\norm{\vmult}
&=\left\lVert m_1v_2+m_2v_1+v_1v_2
 +(v_1a_2+v_2a_1)t
 +\frac{t}{q}\sum_{i=0}^{j_1+j_2}\epsilon_i s^i\right\rVert\\
&\leq \norm{m_1v_2}+\norm{m_2v_1}+\norm{v_1v_2}
 +\norm{(v_1a_2+v_2a_1)t}
 +\frac{t}{q}\left\lVert\sum_{i=0}^{j_1+j_2}\epsilon_i s^i\right\rVert\\
&\leq N_{m_1}\norm{m_1}\norm{v_2}
 +N_{m_2}\norm{m_2}\norm{v_1}
 +n\norm{v_1}\norm{v_2}\\
&\quad+n\norm{v_1}
 \left(\frac{t}{2}\frac{n^{j_2+1}-1}{n-1}
       +\norm{m_2}+\norm{v_2}\right)\\
&\quad+n\norm{v_2}
 \left(\frac{t}{2}\frac{n^{j_1+1}-1}{n-1}
       +\norm{m_1}+\norm{v_1}\right)
 +\frac{t}{2q}\left(\frac{n^{j_1+j_2+1}-1}{n-1}\right)\\
&=
 \left[(N_{m_1}+n)\norm{m_1}
       +\frac{nt}{2}\frac{n^{j_1+1}-1}{n-1}\right]\norm{v_2}\\
&\quad+
 \left[(N_{m_2}+n)\norm{m_2}
       +\frac{nt}{2}\frac{n^{j_2+1}-1}{n-1}\right]\norm{v_1}\\
&\quad+3n\norm{v_1}\norm{v_2}
 +\frac{t}{2q}\left(\frac{n^{j_1+j_2+1}-1}{n-1}\right).
\end{align*}
\hfill$\square$

\section*{Relinearization}

\begin{lemma}[Lemma 12]
Let $\ct$ be a ciphertext of size $M+1$ encrypting $m$, and having
noise $v$. Let $\ctrelin$ of size $N+1$ be the ciphertext encrypting
$m$, obtained by the relinearization of $\ct$, where
$2\leq N+1<M+1$. Then the noise $\vrelin$ in $\ctrelin$ is given by
\[
 \vrelin
 =v-\frac{t}{q}
   \sum_{j=0}^{M-N-1}\sum_{i=0}^{\ell}
   e_{M-j,i}\,c_{M-j}^{(i)},
\]
and can be bounded as
\[
 \norm{\vrelin}\leq \norm{v}
       +\frac{t}{q}(M-N)nB(\ell+1)w.
\]
\end{lemma}

\textbf{Proof.}
Relinearization of a ciphertext from size $M+1$ to size $N+1$,
where $2\leq N+1<M+1$, consists of $M-N$ ``one-step''
relinearizations. In each step, the current ciphertext
$(c_0,c_1,\ldots,c_k)$ is transformed to an intermediate ciphertext
$\ct'=(c'_0,c'_1,\ldots,c'_{k-1})$ using the appropriate evaluation key
\[
 \mathrm{evk}_k=
 \left[
   \left([-(a_{k,i}s+e_{k,i})+w^is^k]_q,\ a_{k,i}\right)
   : i=0,\ldots,\ell
 \right].
\]
In the following step, $\ct'$ becomes the current ciphertext, and so on
until the intermediate ciphertext produced is of size $N+1$, at which
point it is output as $\ctrelin$.

The input ciphertext is $\ct=(c_0,c_1,\ldots,c_M)$, and after the
first one-step relinearization, the intermediate ciphertext is
$\ct'=(c'_0,c'_1,\ldots,c'_{M-1})$, where
\[
 c'_0=c_0+\sum_{i=0}^{\ell}\mathrm{evk}_M[i][0]c_M^{(i)},
 \qquad
 c'_1=c_1+\sum_{i=0}^{\ell}\mathrm{evk}_M[i][1]c_M^{(i)},
\]
and $c'_j=c_j$ for $2\leq j\leq M-1$.
So, for some integer-coefficient polynomials $a_i$, where
$0\leq i\leq \ell+1$,
\begin{align*}
\frac{t}{q}\ct'(s)
&=\frac{t}{q}(c'_0+c'_1s+\cdots+c'_{M-1}s^{M-1})\\
&=\frac{t}{q}\left[
 c_0+\sum_{i=0}^{\ell}\mathrm{evk}_M[i][0]c_M^{(i)}
 +\left(c_1+\sum_{i=0}^{\ell}\mathrm{evk}_M[i][1]c_M^{(i)}\right)s
 +\cdots+c_{M-1}s^{M-1}
 \right]\\
&=\frac{t}{q}\left(
 \sum_{i=0}^{\ell}\mathrm{evk}_M[i][0]c_M^{(i)}
 +s\sum_{i=0}^{\ell}\mathrm{evk}_M[i][1]c_M^{(i)}
 \right)
 +\frac{t}{q}(c_0+c_1s+\cdots+c_{M-1}s^{M-1})\\
&=\frac{t}{q}\left(
 -\sum_{i=0}^{\ell}e_{M,i}c_M^{(i)}
 +\sum_{i=0}^{\ell}a_iqc_M^{(i)}
 +s^M\sum_{i=0}^{\ell}w^ic_M^{(i)}
 \right)
 +\frac{t}{q}(c_0+c_1s+\cdots+c_{M-1}s^{M-1})\\
&=-\frac{t}{q}\sum_{i=0}^{\ell}e_{M,i}c_M^{(i)}
 +\frac{t}{q}\left(c_0+c_1s+\cdots+c_{M-1}s^{M-1}+c_Ms^M\right)
 +t\sum_{i=0}^{\ell}a_ic_M^{(i)}\\
&=m+v-\frac{t}{q}\sum_{i=0}^{\ell}e_{M,i}c_M^{(i)}
 +\left(a_{\ell+1}+\sum_{i=0}^{\ell}a_ic_M^{(i)}\right)t.
\end{align*}
Hence, the noise grows by an additive factor
$-\frac{t}{q}\sum_{i=0}^{\ell}e_{M,i}c_M^{(i)}$
in a one-step relinearization. Iterating this process gives
\[
 \vrelin=v-\frac{t}{q}
 \sum_{j=0}^{M-N-1}\sum_{i=0}^{\ell}
 e_{M-j,i}c_{M-j}^{(i)}.
\]
Bounding $\norm{\vrelin}$,
\begin{align*}
\norm{\vrelin}
&=\left\lVert v-\frac{t}{q}
 \sum_{j=0}^{M-N-1}\sum_{i=0}^{\ell}
 e_{M-j,i}c_{M-j}^{(i)}\right\rVert\\
&\leq \norm{v}+\frac{t}{q}
 \sum_{j=0}^{M-N-1}\sum_{i=0}^{\ell}
 \norm{e_{M-j,i}c_{M-j}^{(i)}}\\
&\leq \norm{v}+\frac{t}{q}(M-N)nB(\ell+1)w.
\end{align*}
\hfill$\square$

\section*{Plain Multiplication}

\begin{lemma}[Lemma 13]
Let $\ct=(x_0,\ldots,x_j)$ be a ciphertext encrypting $m_1$ with
noise $v$, and let $m_2$ be a plaintext polynomial. Let $N_{m_2}$
be an upper bound on the number of non-zero terms in $m_2$.
Let $\ctpmult$ denote the ciphertext obtained by plain multiplication
of $\ct$ with $m_2$. Then the noise in the plain product $\ctpmult$
is $\vpmult=m_2v$, and
\[
  \norm{\vpmult}\leq N_{m_2}\norm{m_2}\norm{v}.
\]
\end{lemma}

\textbf{Proof.}
By definition, $\ctpmult=(m_2x_0,\ldots,m_2x_j)$. Hence, for some
integer-coefficient polynomials $a,a'$,
\begin{align*}
\frac{t}{q}\ctpmult(s)
&=\frac{t}{q}(m_2x_0+m_2x_1s+\cdots+m_2x_js^j)\\
&=m_2\frac{t}{q}(x_0+x_1s+\cdots+x_js^j)\\
&=m_2\frac{t}{q}\ct(s)\\
&=m_2(m_1+v+at)\\
&=m_1m_2+m_2v+m_2at\\
&=[m_1m_2]_t+m_2v+(m_2a-a')t,
\end{align*}
where in the last line we used $[m_1m_2]_t=m_1m_2+a't$.
Hence $\vpmult=m_2v$ and
\[
  \norm{\vpmult}\leq N_{m_2}\norm{m_2}\norm{v}.
\]
\hfill$\square$

\section*{Plain Addition}

\begin{lemma}[Lemma 14]
Let $\ct=(x_0,\ldots,x_j)$ be a ciphertext encrypting $m_1$ with
noise $v$, and let $m_2$ be a plaintext polynomial. Let $\ctpadd$
denote the ciphertext obtained by plain addition of $\ct$ with $m_2$.
Then the noise in $\ctpadd$ is
\[
  \vpadd=v-\frac{r_t(q)}{q}m_2,
\]
and
\[
  \norm{\vpadd}\leq \norm{v}+\frac{r_t(q)}{q}\norm{m_2}.
\]
\end{lemma}

\textbf{Proof.}
By definition of plain addition,
$\ctpadd=(x_0+\Delta m_2,x_1,\ldots,x_j)$. Hence, for some
integer-coefficient polynomials $a,a'$,
\begin{align*}
\frac{t}{q}\ctpadd(s)
&=\frac{t}{q}(x_0+\Delta m_2+x_1s+\cdots+x_js^j)\\
&=\frac{\Delta t}{q}m_2+\frac{t}{q}(x_0+x_1s+\cdots+x_js^j)\\
&=\frac{\Delta t}{q}m_2+\frac{t}{q}\ct(s)\\
&=m_1+v+\frac{q-r_t(q)}{q}m_2+at\\
&=m_1+m_2+v-\frac{r_t(q)}{q}m_2+at\\
&=[m_1+m_2]_t+v-\frac{r_t(q)}{q}m_2+(a-a')t,
\end{align*}
where in the last line we used
$[m_1+m_2]_t=m_1+m_2+a't$. Hence
\[
  \vpadd=v-\frac{r_t(q)}{q}m_2,
\]
and
\[
  \norm{\vpadd}\leq \norm{v}+\frac{r_t(q)}{q}\norm{m_2}.
\]
\hfill$\square$

\section*{Negation}

\begin{lemma}[Lemma 15]
Let $\ct$ be a ciphertext encrypting $m$ with noise $v$, and let
$\ctneg$ be its negation. The noise $\vneg$ in $\ctneg$ is given by
$\vneg=-v$, and
\[
  \norm{\vneg}=\norm{v}.
\]
\end{lemma}

\textbf{Proof.}
If $\ct=(c_0,c_1,\ldots,c_k)$, then its negation is
$\ctneg=(-c_0,-c_1,\ldots,-c_k)=-(c_0,c_1,\ldots,c_k)$.
Thus
\begin{align*}
\frac{t}{q}\ctneg(s)
&=-\frac{t}{q}\ct(s)\\
&=-(m+v+at)\\
&=-m+(-v)+(-a)t.
\end{align*}
Hence the noise $\vneg$ in $\ctneg$ is $-v$ and
$\norm{\vneg}=\norm{v}$.
\hfill$\square$

\section*{Subtraction}

Suppose $\ct_1$ and $\ct_2$ are two ciphertexts encrypting $m_1$ and
$m_2$, and we want to compute a ciphertext $\ctsub$ encrypting
$m_1-m_2$. We could first negate $\ct_2$ to obtain a ciphertext
$\ct'_2$ that encrypts $-m_2$, and then perform an addition of
$\ct_1$ and $\ct'_2$. By viewing subtraction in this way, the noise
growth in subtraction is at most that for addition, since the noise
does not change in norm in negation.

\begin{lemma}[Lemma 16]
Let $\ct_1$ and $\ct_2$ be two ciphertexts encrypting $m_1,m_2$,
respectively, with noises $v_1,v_2$, respectively. The noise
$\vsub$ in the result $\ctsub$ is bounded by
\[
  \norm{\vsub}\leq \norm{v_1}+\norm{v_2}.
\]
\end{lemma}

\section*{Plain Subtraction}

By the same argument as for subtraction, the noise growth in plain
subtraction is at most that for plain addition.

\begin{lemma}[Lemma 17]
Let $\ct$ be a ciphertext encrypting $m_1$ with noise $v$, and let
$m_2$ be a plaintext polynomial. Let $\ctpsub$ denote the ciphertext
obtained by plain subtraction of $m_2$ from $\ct$. Then the noise
$\vpsub$ in $\ctpsub$ is bounded by
\[
  \norm{\vpsub}\leq \norm{v}+\frac{r_t(q)}{q}\norm{m_2}.
\]
\end{lemma}

\section*{Transcription Verification}

The following checks were performed against the rendered PDF pages
27--34 of the supplied SEAL 2.3.1 manual:

\begin{enumerate}[leftmargin=2em]
\item Lemma numbers and section order match the appendix: Lemmas 9--17
      (Initial Noise, Addition, Multiplication, Relinearization,
      Plain Multiplication, Plain Addition, Negation, Subtraction,
      Plain Subtraction).
\item The displayed bound in Lemma 9 matches
      $\frac{r_t(q)}{q}\norm{m}+\frac{tB}{q}(2n+1)$.
\item The three terms multiplying $\norm{v_2}$, $\norm{v_1}$, and
      $\norm{v_1}\norm{v_2}$, plus the final rounding term in Lemma 11,
      were checked against the printed statement and the last line
      of its proof. In particular, the powers are $n^{j_1+1}$,
      $n^{j_2+1}$, and $n^{j_1+j_2+1}$.
\item The relinearization sum in Lemma 12 uses indices
      $j=0,\ldots,M-N-1$ and $i=0,\ldots,\ell$, with
      $e_{M-j,i}c_{M-j}^{(i)}$; its bound is
      $\norm{v}+\frac{t}{q}(M-N)nB(\ell+1)w$.
\item Lemmas 13--17 were checked against their displayed statements
      and, where present, their proofs. Lemma 16 is a stated bound
      supported by the preceding explanatory paragraph; Lemma 17
      has no separate proof in the manual.
\end{enumerate}

\textbf{Transcription status:} the mathematical expressions above
match the corresponding appendix statements in the supplied PDF.
The original source's notation is retained; typographic layout and
line breaks have been normalized for readable LaTeX.

\end{document}
